
% <<<PDFPAGE 43>>>
\chapter{推理时计算扩展}

\section{新的扩展范式}

多年来，占主导地位的扩展定律（scaling law）认为，训练计算量越大，模型就越好。Chinchilla 的最优训练方案建议将计算预算平均分配给模型规模和训练词元数。2024 至 2025 年出现了一种补充范式，称为\textbf{推理时计算扩展（inference-time compute scaling）}，即增加推理阶段的计算投入来提高输出质量。

在许多任务中，尤其在推理、数学和编程任务上，让小模型“多思考一会儿”，可能比让大模型“立即作答”表现更好。问题的重心由“我们能训练多大的模型”转向“模型运行时，我们能为推理过程投入多少计算”。

\section{思维链及其推理成本}

思维链（Chain-of-Thought，CoT）提示（Wei 等，2022）要求模型在给出最终答案前先生成中间推理步骤。这显著改善了多步推理任务上的表现，但也需要额外计算。一个答案只需 10 个词元的问题，推理过程可能需要 500 个词元，使推理成本变为原来的 50 倍。

\textbf{推理时计算扩展所问的问题。}给定固定的计算预算，是运行一次较大的模型更好，还是让较小的模型生成更多推理词元（reasoning tokens）更好？

近期研究（Snell 等，2024，《Scaling LLM Test-Time Compute Optimally》）表明，答案取决于问题的难度。对于难题，让小模型进行更多推理通常更高效。对于简单问题，延长推理过程则会浪费计算。

\section{过程奖励模型与搜索}

\textbf{过程奖励模型（Process Reward Model，PRM）}通过训练来评估各个中间推理步骤的质量，评价范围不限于最终答案。有了 PRM，就可以开展搜索。先生成多条候选推理路径，再用 PRM 选择最有希望的路径或对候选路径重新排序。

常见搜索策略包括以下三种。

\begin{itemize}
  \item \textbf{\(N\) 选一（Best-of-\(N\)）。}生成 \(N\) 个独立的解答，选出 PRM 得分最高的一个，也可以用多数投票选出答案。
  \item \textbf{束搜索（beam search）。}保留 \(K\) 条未完成的推理路径，每一步都扩展最有希望的路径。
  \item \textbf{蒙特卡洛树搜索（Monte Carlo Tree Search，MCTS）。}在树搜索中用 PRM 充当价值函数。
\end{itemize}

当 \(N=64\) 时，\(N\) 选一在数学评测（MATH、AIME）上可以达到或超过规模为其 10 倍的模型。

% <<<PDFPAGE 44>>>
\section{DeepSeek-R1 与 o1 式架构}

OpenAI o1 和 DeepSeek-R1 代表了一个质的转变。与用提示引出思维链的方法相比，这类模型通过强化学习（Reinforcement Learning，RL）习得推理能力。模型在给出最终答案前会生成长段内部独白（“思考”）。原书将其训练描述为通过基于人类反馈的强化学习（Reinforcement Learning from Human Feedback，RLHF）或群组相对策略优化（Group Relative Policy Optimization，GRPO），根据可验证的奖励信号训练这些思考过程，奖励信号例如数学题答案是否正确。原书写出的目标为
\[
  \mathcal{L}_{\mathrm{GRPO}}
  = -\mathbb{E}_o\left[
    \min\left(
      \frac{\pi_\theta(o\mid q)}{\pi_{\mathrm{ref}}(o\mid q)} A_o,
      \operatorname{clip}(\cdot,1\pm\epsilon) A_o
    \right)
  \right]+\beta D_{\mathrm{KL}}(\pi_\theta\Vert\pi_{\mathrm{ref}}).
\]
\translatornote{上式照原书保留。DeepSeek-R1 论文对同一问题的一组输出取平均，概率比的分母为旧策略 \(\pi_{\theta_{\mathrm{old}}}\)，参考策略 \(\pi_{\mathrm{ref}}\) 用于 KL 散度项。参见 \href{https://arxiv.org/html/2501.12948v1}{论文式（1）}。}

这样的模型处理难题时可用数千个词元来“思考”，以增加的推理成本换取准确率。DeepSeek-R1-Zero 表明，即便没有监督式推理轨迹（reasoning traces），这种推理行为也能通过纯强化学习涌现。

\begin{interview}
\question{1}{什么是推理时计算扩展？它与训练时计算扩展有何区别？}
训练时计算扩展（training-time compute scaling）（Kaplan 等，2020，Hoffmann 等，2022，Chinchilla）研究模型能力怎样随训练浮点运算预算增长。参数更多、训练词元更多，模型表现就更好。推理时计算扩展是一个较新的范式，研究给定一个已训练好的固定模型时，增加推理计算量会怎样提高性能。

可以通过以下方式投入更多推理计算。（1）用思维链生成更多推理词元。（2）生成多个候选解答并从中选出最佳解答，即 \(N\) 选一采样。（3）用验证器（verifier）或奖励模型进行树搜索。（4）进行迭代改进（iterative refinement），由模型评议并改善自己的输出。

额外的推理计算可以替代一部分训练计算。例如，在某些推理任务上，70 亿参数模型投入 64 倍的推理计算，可以达到 700 亿参数模型的表现。这改变了 AI 部署的成本结构，计算预算从一次性的训练支出转移到按查询计的推理支出。

\question{2}{请解释过程奖励模型。它与结果奖励模型有什么区别？}
\textbf{结果奖励模型（Outcome Reward Model，ORM）}仅根据最终答案是否正确，给完整解答打一个奖励分。即使两种解法都得出了正确答案，ORM 也无法分辨其中哪一种推理有误、哪一种推理成立。ORM 无法引导对中间步骤的搜索。

\textbf{过程奖励模型}经过训练，可以为每个推理步骤的正确性打分，因而能在整条推理链上提供密集反馈。PRM 需要在训练数据中有逐步标注。尽管已有自动化方法，收集标注仍费人力。推理时，PRM 支持逐步搜索，因而可以进行蒙特卡洛树搜索、束搜索和逐步重排序。在重视中间路径的复杂多步推理中，PRM 的表现优于 ORM，可以避免模型因错误推理碰巧得到正确答案。训练和应用 PRM 的成本则显著更高。

\question{3}{什么是 \(N\) 选一采样？它在什么情况下高效，什么情况下低效？}
% <<<PDFPAGE 45>>>
\(N\) 选一（BoN）采样在温度 \(>0\) 的条件下，让模型生成 \(N\) 个独立的候选输出，再由奖励模型或验证器逐个打分，返回得分最高的输出。这是最简单的推理时计算扩展方法之一。

\textbf{高效的情形。}（1）任务有可验证的正确答案，如数学、编程和形式逻辑，且验证器廉价可靠。（2）生成结果多样。如果 \(N\) 个样本都相同，BoN 就无益处。（3）\(N\) 较小（\(\le 32\)），且每个样本的推理成本较低。

\textbf{低效的情形。}（1）开放式任务没有可靠的评估器，如创意写作，难以自动评价主观质量。（2）\(N\) 极大。在许多情形中，BoN 的质量按 \(O(\log N)\) 的量级增长，对应于从一个分布中抽取 \(N\) 个样本后最大值的期望，因而边际收益递减。（3）模型本身推理成本很高，生成 \(N\) 个样本的成本是单个样本的 \(N\) 倍。同样的计算预算可以更有效地用于束搜索或逐步改进。

\question{4}{请描述 DeepSeek-R1 采用的 GRPO 训练目标。它和 PPO 有什么区别？}
GRPO 是一种用于微调大语言模型的强化学习算法，无须单独的价值网络（value network，也称 critic）。近端策略优化（Proximal Policy Optimization，PPO）用价值网络估计优势函数（advantage function）\(A\)。GRPO 则针对同一个问题生成一组输出，计算每个输出相对于组内其他输出的优势，
\[
  A_i=(r_i-\operatorname{mean}(\mathbf{r}))/\operatorname{std}(\mathbf{r}),
\]
其中 \(\mathbf{r}\) 是该组的奖励向量。这样无须训练价值函数，可以节省内存和计算。组内归一化会降低方差，也避免了估计价值函数时产生的不稳定问题。

与 PPO 相比，GRPO 更简单，所需的超参数更少，而且省去了价值网络，因而更节省内存。不过，由于基准值（baseline）是根据当前批次估计的，单次更新的方差可能更大。DeepSeek-R1-Zero 表明，仅靠 GRPO 和简单的可验证奖励，就能诱发自我验证、长时间思考等复杂推理行为。

\question{5}{如果要设计数学辅导应用，你会使用更大的基座模型，还是选较小的模型并增加推理时的计算量？请说明理由。}
我会选\textbf{较小的模型，并增加推理时的计算量}，理由有四点。

（1）\textbf{数学答案可以验证。}正确答案可以用程序检查，因此可以采用可靠的 \(N\) 选一或 PRM 引导的搜索，无须人工评估。

（2）\textbf{推理轨迹具有教学价值。}能分步展示解题过程的数学辅导工具比只给出答案的工具更有用。推理时生成思维链本身就是教学功能。

（3）\textbf{单次查询成本较低。}即使计入额外生成的词元，部署一个增加推理计算的 70 亿参数 R1 式模型，每次查询的成本也低于部署一个 700 亿参数的稠密模型。

（4）\textbf{计算预算可以调整。}简单问题可以少用一些推理词元，难题可以多用一些。与固定的大模型相比，这样能随题目调节计算投入。

实际部署时，我会选 70 亿或 140 亿参数的微调推理模型，如 DeepSeek-R1-Distill 或同类模型，并让推理预算可配置，以此替代 700 亿参数的基准模型。
\end{interview}
